% PS5 mathematical companion. Build with make in this directory.
% The slide order follows the assignment README; no reference results appear here.
% The starter docstrings cite four slides by number and title: 4 'Pricing the
% monthly call', 5 'One month of the covered call fund', 6 'The growth rate over
% a roll', and 11 'Advanced: one month of the collar fund'. Update those
% citations if the slide order changes.
\documentclass[aspectratio=169,11pt]{beamer}
\usepackage{vnslides}
\newcommand{\Prem}{\mathcal{P}}
\newcommand{\Dfac}{\mathcal{D}}
\newcommand{\one}{\mathbf{1}}

\title{\texorpdfstring{{\vnheavy PS5}: Does a Covered Call\\Fund Beat Its Shares?}{PS5: Does a Covered Call Fund Beat Its Shares?}}
\subtitle{CHEME 4/5660 Cornell University | Mathematical Companion}
\author{Copyright Jeffrey Varner 2026. All Rights Reserved}
\begin{document}
\vntitlepage

\begin{frame}{Disclaimer and risks}
\textbf{Training only.} This assignment uses historical option quotes, a
simple price model, and simplified trading assumptions for teaching. It
provides no investment advice.

\vspace{1em}
\textbf{Covered calls.} A covered call fund keeps the full loss of its shares,
less the premiums, and gives up every gain above each month's strike. A collar
limits each month's loss with a put, at the cost of that put's premium.

\vspace{1em}
\textbf{Simplifying assumptions.} One call (and, for the collar, one put) per
share, measured per share. Every option is held to expiration and settled at
the close that day, with no early exercise. We omit commissions, dividends,
and taxes. The quotes come from an indicative data feed.
\end{frame}

\begin{frame}{The central question}
Does a covered call fund beat simply holding the shares, and what does it
give up?
\vspace{0.2em}
\begin{itemize}\setlength{\itemsep}{0.15em}
\item \hilite{Study 1: the promise.} Simulate six months of share prices for
each ticker, price the calls on the binomial tree, and compare the
distributions of the fund's and the shares' growth rates.
\item \hilite{Study 2: what happened.} Replay each ticker's six real monthly
calls at their market prices, and compare the fund with the shares, SPY, and
cash over the same days.
\item \hilite{Study 3: the conclusion.} Decide what the fund gave up and what
it got, and what one half-year of history can tell us.
\end{itemize}
As in the course example, $S_j$ is the share price at the start of month $j$,
$K_j$ the strike of the call sold then, and $\Prem_{c,j}$ its premium (USD/share).
\slidenote{Reading guide}{Both tracks do Studies 1--3. The collar fund on slide \pageref{fr:collar} is Advanced only.}
\end{frame}

\begin{frame}{Pricing the monthly call}
The call sold at the start of month $j$, strike $K_j$, expires after $\Delta t$
years. We price it on the Cox--Ross--Rubinstein (CRR) binomial tree with $h=60$
steps of length $\tau=\Delta t/h$:
\[
 u=e^{\sigma\sqrt{\tau}},\qquad d=\frac{1}{u},\qquad
 \bar p=\frac{e^{g_y\tau}-d}{u-d},
\]
where $\sigma$ is the pricing volatility and $g_y$ the continuously
compounded risk-free rate. Rolling the payoff $(S-K_j)^{+}$ back through the
tree under $\bar p$, discounting by $e^{-g_y\tau}$ per step and keeping at
each node the larger of that value and the exercise value (American), gives
$\Prem_{c,j}$.

In code: \texttt{build(MyAmericanCallContractModel, ...)} with strike $K_j$,
then \texttt{binomial\_premium(call; method = :crr, h = 60, ...)} with time
$\Delta t$, volatility $\sigma$, share price $S_j$, and rate $g_y$.
\slidenote{No dividends}{An American call on a stock without dividends is never worth exercising early, so it is priced like a European call.}
\end{frame}

\begin{frame}{One month of the covered call fund}
The fund's wealth $W_j$ is its shares at $S_j$. It sells one call per share and
holds the premium as cash, which grows by $\Dfac_{T,0}(g_y)=e^{g_y T}$ over the
call's life $T$. At expiration the shares are worth $S_{j+1}$, less the call's payoff:
\[
 \frac{W_{j+1}}{W_j}
 =\frac{S_{j+1}-\left(S_{j+1}-K_j\right)^{+}+\Dfac_{T,0}(g_y)\cdot\Prem_{c,j}}{S_j}
 =\frac{\min\left(S_{j+1},K_j\right)+\Dfac_{T,0}(g_y)\cdot\Prem_{c,j}}{S_j}.
\]
Holding the shares instead gives $S_{j+1}/S_j$.
\begin{itemize}\setlength{\itemsep}{0.1em}
\item \hilite{Expires worthless,} $S_{j+1}\leq K_j$: the fund leads the
shares by $\Dfac_{T,0}(g_y)\cdot\Prem_{c,j}/S_j$.
\item \hilite{Exercised,} $S_{j+1}>K_j$: the ratio stops at
$\bigl(K_j+\Dfac_{T,0}(g_y)\cdot\Prem_{c,j}\bigr)/S_j$.
\end{itemize}
\slidenote{Test month}{$S_j=100$, $K_j=102$, $\Prem_{c,j}=0.87$, $\Dfac=1.00345$: 1.0053 at $S_{j+1}=99.66$; 1.0287 for any $S_{j+1}\geq102$.}
\end{frame}

\begin{frame}{The growth rate over a roll}
Over $N$ months that last $T_{\text{roll}}$ years in all, the fund's growth rate
is the log of its final over its initial wealth, divided by the time. The
final wealth is the product of the monthly ratios, so:
\[
 g=\frac{1}{T_{\text{roll}}}\ln\left(\frac{W_{N+1}}{W_1}\right)
 =\frac{1}{T_{\text{roll}}}\sum_{j=1}^{N}\ln\left(\frac{W_{j+1}}{W_j}\right)
 \qquad(\text{units: }1/\text{yr}).
\]
The same formula with $S_{j+1}/S_j$ gives the shares' growth rate. The order of
the ratios does not matter, and a ratio of exactly 1 adds nothing.

\vspace{0.5em}
\begin{itemize}\setlength{\itemsep}{0.3em}
\item \hilite{In the model:} $N=6$ months of $\Delta t=D_1/365$ year, the
first trade's term, so $T_{\text{roll}}=6\,\Delta t$.
\item \hilite{In history:} $T_{\text{roll}}$ is the calendar days from a
ticker's first entry to its last expiration, divided by 365, about $0.47$ year.
\end{itemize}
\end{frame}

\begin{frame}{The model of the share price}
For each ticker, the model simulates the share price at the start of each
month with geometric Brownian motion (GBM), our belief about the share price:
\[
 S_{j+1}=S_j\exp\Bigl[\bigl(\hat\mu-\tfrac{\hat\sigma^{2}}{2}\bigr)\Delta t
 +\hat\sigma\sqrt{\Delta t}\,Z_j\Bigr],\qquad Z_j\sim\mathcal{N}(0,1),
\]
on 10,000 paths from $S_1=S(0)$, the close on the first entry date, with
$\Delta t=D_1/365$, the first trade's term. The drift
$\hat\mu=\hat g+\hat\sigma^{2}/2$ and volatility $\hat\sigma$ come from the
252 daily changes before that date ($\hat g$: mean log growth rate).

\vspace{0.3em}
The strike is $K_j=m\,S_j$, with $m=K/S(0)$ from the first trade. The call is
priced at one of two volatilities:
\begin{itemize}\setlength{\itemsep}{0.1em}
\item \hilite{historical,} $\sigma=\hat\sigma$, as in the course example;
\item \hilite{implied,} $\sigma=\sigma_{\text{IV}}$, implied by the call's
market price on the first entry date.
\end{itemize}
\slidenote{Two probabilities}{Paths use our belief (drift $\hat\mu$); premiums use the risk-neutral tree (drift $g_y$).}
\end{frame}

\begin{frame}{The premium is a fixed fraction of the share price}
Every node of the tree is $S_j$ times a factor that does not depend on $S_j$,
and the payoff $(S-mS_j)^{+}$ scales the same way. So the premium is
proportional to the share price:
\[
 \Prem_{c,j}=c(m,\sigma)\,S_j,\qquad
 c(m,\sigma)=\frac{\Prem_c\bigl(S(0),\,m\,S(0)\bigr)}{S(0)}.
\]
The report prices one call per ticker and pricing volatility, and the monthly
ratio depends only on that month's price change:
\[
 \frac{W_{j+1}}{W_j}=\min\left(\frac{S_{j+1}}{S_j},\,m\right)
 +\Dfac_{T,0}(g_y)\cdot c(m,\sigma).
\]
The fund gives up every monthly price ratio above $m$, and collects
$c(m,\sigma)$ every month. A farther strike (larger $m$) gives up less and
collects less.
\slidenote{Report}{\textit{Premiums: model and market} lists $c(m,\hat\sigma)$, $c(m,\sigma_{\text{IV}})$, and the market's mid price over $S(0)$.}
\end{frame}

\begin{frame}{Replaying history}
Each ticker has six back-to-back trades. Trade $j$ enters at the close
$S_{0,j}$, sells the call with strike $K_j$ at its mid price $\Prem_{c,j}$, and
settles at the close $S_{T,j}$ on the expiration date, $D_j$ calendar days
later:
\[
 \frac{W_{j+1}}{W_j}
 =\frac{\min\left(S_{T,j},K_j\right)+\Dfac_j\cdot\Prem_{c,j}}{S_{0,j}},
 \qquad \Dfac_j=e^{g_y D_j/365}.
\]
When trade $j+1$ enters after trade $j$ expires, the fund holds only its shares
in between, so the fund and the shares both gain the ratio
$S_{0,j+1}/S_{T,j}$. The risk-free rate comes from the 13-week bill yield
$y=0.0376$ and term $\tau=91/365$, as in the course example:
\[
 g_y=\frac{1}{\tau}\ln\left(1+y\,\tau\right)\approx0.0374\ \text{per year}.
\]
\slidenote{Benchmarks}{SPY bought at its close on the ticker's first entry date and sold on its last expiration date; cash grows at $g_y$, so its growth rate is $g_y$.}
\end{frame}

\begin{frame}{Promise versus outcome}
Here the simulation follows ticker $i$'s real roll: each trade keeps its real
$K_j/S_{0,j}$, $\Prem_{c,j}/S_{0,j}$, and $D_j$, with the gaps; only prices are
simulated. It gives $N_p=10{,}000$ pairs
$\bigl(g_{\text{fund}}^{(p)},g_{\text{shares}}^{(p)}\bigr)$, and with the
indicator $\one[\cdot]$:
\[
 \hat P_i=\frac{1}{N_p}\sum_{p=1}^{N_p}
 \one\Bigl[g_{\text{fund}}^{(p)}>g_{\text{shares}}^{(p)}\Bigr],
 \qquad
 \text{expected count}=\sum_{i=1}^{29}\hat P_i .
\]
For one ticker, with the real difference $\delta_i=g_{\text{fund}}-g_{\text{shares}}$,
the model percentile is:
\[
 q_i=\frac{1}{N_p}\sum_{p=1}^{N_p}
 \one\Bigl[g_{\text{fund}}^{(p)}-g_{\text{shares}}^{(p)}\leq\delta_i\Bigr].
\]
\slidenote{Reading}{A percentile near 0\% or 100\% is rare under the model; with matched trades it points at the share prices.}
\end{frame}

\begin{frame}{Advanced: one month of the collar fund}\label{fr:collar}
The collar fund also buys a put with strike $K_1<K_2$, where $K_2$ is the
call's strike. With the put's premium $\Prem_1$ and the call's $\Prem_2$, the
month's payoffs add to the share price held between a floor and a cap:
\[
 S_{j+1}+\left(K_1-S_{j+1}\right)^{+}-\left(S_{j+1}-K_2\right)^{+}
 =\min\Bigl(\max\left(S_{j+1},K_1\right),\,K_2\Bigr).
\]
The net premium $\Prem_2-\Prem_1$, negative when the put costs more, grows to
expiration with the cash:
\[
 \frac{W_{j+1}}{W_j}
 =\frac{\min\bigl(\max\left(S_{j+1},K_1\right),K_2\bigr)
 +\Dfac_{T,0}(g_y)\cdot\left(\Prem_2-\Prem_1\right)}{S_j}.
\]
The put is the trade's \texttt{long\_put} leg, near 0.10 delta, well below the
share price. The model uses $K_1=m_1S_j$ and $K_2=m\,S_j$, with $m_1$ and $m$ the
first trade's strikes over $S(0)$, and prices the put at its implied volatility.
\slidenote{Course example}{\textit{A Protective Collar on SPY} uses the same floor $K_1$, cap $K_2$, and premiums $\Prem_1$ (put) and $\Prem_2$ (call).}
\end{frame}

\begin{frame}{Summary}
PS5 runs a covered call fund twice: in a model of the share price, and on the
real monthly calls of 29 tickers, and compares it with holding the shares.

\vspace{0.6em}
\textbf{Key takeaways}
\begin{itemize}\setlength{\itemsep}{0.5em}
\item Each month the fund's wealth ratio is the share price capped at the
strike plus the grown premium, over the starting price. The premiums lift the
falling months and the strike caps the rising ones.
\item The growth rate over a roll is the sum of the logs of the monthly wealth
ratios divided by the time in years, in 1/yr for the model and for history.
\item A model's probability that the fund beats the shares is a promise about
many possible half-years; history gives one half-year, shared by all 29
tickers.
\end{itemize}
\vspace{0.6em}
Use these measures to answer the central question.
\end{frame}
\end{document}
